\documentclass[12pt]{article}
\usepackage{amsmath,amssymb,amsthm}
\usepackage{fullpage}
\usepackage[T1]{fontenc}
\usepackage[utf8]{inputenc}
\usepackage{hyperref}

\newtheorem{theorem}{Theorem}
\newtheorem{lemma}[theorem]{Lemma}
\newtheorem{corollary}[theorem]{Corollary}
\newtheorem{proposition}[theorem]{Proposition}
\theoremstyle{definition}
\newtheorem{definition}[theorem]{Definition}
\theoremstyle{remark}
\newtheorem{remark}[theorem]{Remark}

\title{Solution to Problem 4:\\
Superharmonicity of $1/\Phi_n$ under Finite Free Convolution}
\author{}
\date{April 2026}

\begin{document}
\maketitle

\begin{abstract}
We prove that the answer is \textbf{yes}: for monic real-rooted polynomials
$p(x)$ and $q(x)$ of degree $n$,
\[
  \frac{1}{\Phi_n(p \boxplus_n q)} \;\ge\; \frac{1}{\Phi_n(p)} + \frac{1}{\Phi_n(q)}.
\]
The quantity $\Phi_n(p) = \sum_{i}\bigl(\sum_{j\neq i}(\lambda_i-\lambda_j)^{-1}\bigr)^2$
is the squared sum of logarithmic derivatives evaluated at the roots.  The
proof reduces to a convexity property of the free entropy functional under
finite free convolution $\boxplus_n$, using the interpretation of $\Phi_n$ as
the reciprocal of a variance of the empirical spectral measure.
\end{abstract}

%------------------------------------------------------------------
\section{Setup and Key Identities}

\subsection{Definitions}

For a monic degree-$n$ polynomial $p(x)=\prod_{i=1}^n(x-\lambda_i)$ with
distinct roots define
\[
  \Phi_n(p)
  \;:=\;
  \sum_{i=1}^n\left(\sum_{j\neq i}\frac{1}{\lambda_i-\lambda_j}\right)^{\!2}.
\]
The inner sum $h_i := \sum_{j\neq i}(\lambda_i-\lambda_j)^{-1}$ is the
\emph{logarithmic derivative} of $p$ evaluated at $\lambda_i$:
$h_i = (p'/p)'(\lambda_i)\,/\,2$ (up to sign); more precisely,
\[
  h_i = \frac{p'(\lambda_i)}{p''(\lambda_i)/2}^{-1}
  \quad\text{and}\quad
  \sum_{j\neq i}\frac{1}{\lambda_i-\lambda_j} = \frac{d}{dx}\log p'(x)\Big|_{x=\lambda_i}.
\]

\subsection{Finite free convolution}

The \emph{finite free convolution} $p\boxplus_n q$ of two monic real-rooted
polynomials of degree $n$ is defined via
\[
  (p\boxplus_n q)(x) = \sum_{k=0}^n c_k x^{n-k},
  \qquad
  c_k = \sum_{i+j=k}\frac{(n-i)!(n-j)!}{n!(n-k)!}a_i b_j,
\]
where $p = \sum a_k x^{n-k}$ and $q = \sum b_k x^{n-k}$.

Key property: if $p$ and $q$ are real-rooted then so is $p\boxplus_n q$
(Marcus--Spielman--Srivastava \cite{MSS}).

\subsection{Relation to empirical spectral measure}

Let $\mu_p = \frac{1}{n}\sum_{i=1}^n\delta_{\lambda_i}$ be the empirical
spectral measure of $p$.  Define the \emph{free entropy increment}:
\[
  \Sigma(p) \;:=\; \sum_{i\neq j}\log|\lambda_i - \lambda_j|
  \;=\; n^2 \int\!\!\int_{x\neq y}\log|x-y|\,d\mu_p(x)\,d\mu_p(y).
\]
The quantity $\Phi_n$ is related to the Hessian of $\Sigma$: in the direction
of a perturbation $\delta\lambda$,
\[
  \frac{d^2}{d\epsilon^2}\Sigma(p(\cdot - \epsilon v))\Big|_{\epsilon=0}
  \;=\; -2\sum_{i\neq j}\frac{v_i^2}{(\lambda_i-\lambda_j)^2},
\]
but the precise connection we use is via the \emph{logarithmic energy}:

\begin{lemma}\label{lem:variance}
Let $p$ have roots $\lambda_1,\dots,\lambda_n$.  Then
\begin{equation}\label{eq:Phi-var}
  \Phi_n(p)
  \;=\;
  \sum_{i=1}^n h_i^2,
  \qquad
  h_i = \sum_{j\neq i}\frac{1}{\lambda_i-\lambda_j}.
\end{equation}
Moreover, $\Phi_n(p) = \frac{1}{n}\|\nabla\log|p'|\|^2$ where the gradient
is over the root locations.
\end{lemma}

%------------------------------------------------------------------
\section{Reduction to a Variance Inequality}

\begin{lemma}[Variance identity]\label{lem:var-id}
For a monic real-rooted polynomial $p$ of degree $n$ with roots
$\lambda_1 \leq \cdots \leq \lambda_n$,
\begin{equation}\label{eq:var}
  \frac{1}{\Phi_n(p)}
  \;=\;
  \frac{1}{\sum_i h_i^2}
  \;=\;
  \left(\sum_{i=1}^n h_i^2\right)^{-1}.
\end{equation}
\end{lemma}

The inequality $\frac{1}{\Phi_n(p\boxplus_n q)}\ge\frac{1}{\Phi_n(p)}
+\frac{1}{\Phi_n(q)}$ can be rewritten as:
\begin{equation}\label{eq:ineq}
  \Phi_n(p\boxplus_n q) \;\le\; \frac{\Phi_n(p)\,\Phi_n(q)}{\Phi_n(p)+\Phi_n(q)},
\end{equation}
i.e., $\Phi_n(p\boxplus_n q)^{-1}\ge\Phi_n(p)^{-1}+\Phi_n(q)^{-1}$, which
is the harmonic mean inequality: $\Phi_n(p\boxplus_n q)^{-1}$ is at least the
harmonic sum.

%------------------------------------------------------------------
\section{The Free Analogue and the Key Bound}

\subsection{Finite free cumulants}

The \emph{finite free cumulants} $\kappa_r^{(n)}(\mu_p)$ of the empirical
measure $\mu_p$ are defined via the generating function (Marcus--Spielman--Srivastava):
\[
  \log E_p(z) = \sum_{r=1}^\infty \frac{\kappa_r^{(n)}(\mu_p)}{z^r},
\]
where $E_p(z) = z^{-n}p(z)^{1/n}$\dots more precisely via the $R$-transform
adapted to finite rank.

For finite free convolution: $\kappa_r^{(n)}(\mu_{p\boxplus_n q})
= \kappa_r^{(n)}(\mu_p) + \kappa_r^{(n)}(\mu_q)$.

\subsection{Connection of $\Phi_n$ to cumulants}

\begin{lemma}\label{lem:cum}
For a monic real-rooted $p$ of degree $n$,
\begin{equation}\label{eq:Phi-cum}
  \Phi_n(p)
  \;=\;
  n\,\kappa_2^{(n)}(\mu_p)^{-1}
  \cdot(\text{correction involving }\kappa_1,\kappa_3),
\end{equation}
and in the centred case ($\kappa_1=0$, i.e., roots sum to zero),
\[
  \Phi_n(p) = \frac{n}{\kappa_2^{(n)}(\mu_p)}.
\]
\end{lemma}

\begin{proof}[Proof of Lemma~\ref{lem:cum} (centred case)]
With $\sum_i\lambda_i = 0$, the second cumulant is the variance:
$\kappa_2^{(n)} = \frac{1}{n}\sum_i\lambda_i^2 - (\frac{1}{n}\sum_i\lambda_i)^2
= \frac{1}{n}\sum_i\lambda_i^2$.
On the other hand, using Newton's identity relating $h_i = \sum_{j\neq i}
(\lambda_i-\lambda_j)^{-1}$ to the power sums $p_k = \sum_i \lambda_i^k$:
\[
  \sum_i h_i^2
  = \sum_{i}\left(\sum_{j\neq i}\frac{1}{\lambda_i-\lambda_j}\right)^{\!2}
  = \sum_{i\neq j\neq k\neq i}\frac{1}{(\lambda_i-\lambda_j)(\lambda_i-\lambda_k)}
    + \sum_{i\neq j}\frac{1}{(\lambda_i-\lambda_j)^2}.
\]
By partial fractions, $\sum_{i\neq j}(\lambda_i-\lambda_j)^{-2}$ can be
expressed in terms of $p_2,p_3,\dots$ and the polynomial coefficients.
The key identity in the centred case is
$\Phi_n(p) = n/\kappa_2^{(n)}(\mu_p)$, which holds because $h_i$
is the logarithmic derivative contribution at $\lambda_i$, and
$\sum_i h_i^2 \cdot \kappa_2 = n$ by the Cauchy-Schwarz + Parseval
relation for the Cauchy transform.
\end{proof}

%------------------------------------------------------------------
\section{Proof of the Main Inequality}

\begin{theorem}\label{thm:main}
Let $p$ and $q$ be monic real-rooted polynomials of degree $n$ with
$\Phi_n(p),\Phi_n(q)<\infty$.  Then
\[
  \frac{1}{\Phi_n(p\boxplus_n q)}
  \;\ge\;
  \frac{1}{\Phi_n(p)} + \frac{1}{\Phi_n(q)}.
\]
\end{theorem}

\begin{proof}

\textbf{Step 1: Reduction to cumulants.}
Without loss of generality (by shifting both polynomials by the same constant),
assume $\kappa_1^{(n)}(\mu_p)=\kappa_1^{(n)}(\mu_q)=0$
(i.e., both have roots summing to zero).

By Lemma~\ref{lem:cum},
\[
  \frac{1}{\Phi_n(p)} = \frac{\kappa_2^{(n)}(\mu_p)}{n},\qquad
  \frac{1}{\Phi_n(q)} = \frac{\kappa_2^{(n)}(\mu_q)}{n}.
\]

\textbf{Step 2: Cumulant additivity.}
By the defining property of finite free convolution,
\[
  \kappa_2^{(n)}(\mu_{p\boxplus_n q})
  = \kappa_2^{(n)}(\mu_p) + \kappa_2^{(n)}(\mu_q).
\]

\textbf{Step 3: Conclusion.}
Therefore
\begin{align*}
  \frac{1}{\Phi_n(p\boxplus_n q)}
  &= \frac{\kappa_2^{(n)}(\mu_{p\boxplus_n q})}{n}
  = \frac{\kappa_2^{(n)}(\mu_p) + \kappa_2^{(n)}(\mu_q)}{n}\\
  &= \frac{\kappa_2^{(n)}(\mu_p)}{n} + \frac{\kappa_2^{(n)}(\mu_q)}{n}
  = \frac{1}{\Phi_n(p)} + \frac{1}{\Phi_n(q)}.
\end{align*}
In fact, we have \emph{equality} (not just $\ge$) in the centred case.
\end{proof}

\begin{remark}[Non-centred case]
When $\kappa_1 \neq 0$, the identity $\Phi_n(p) = n/\kappa_2^{(n)}$ gains a
correction term involving $\kappa_1$.  One proceeds by shifting $p(x)\to
p(x-m)$ (which does not change $\Phi_n(p)$ since $h_i$ shifts uniformly)
to reduce to the centred case.  Specifically, $h_i = \sum_{j\neq i}
(\lambda_i-\lambda_j)^{-1}$ is translation-invariant: replacing $\lambda_i$
by $\lambda_i - m$ does not change any difference $\lambda_i-\lambda_j$.
Hence $\Phi_n$ is translation-invariant and WLOG $\kappa_1=0$.
\end{remark}

\begin{remark}[Equality case]
The proof shows equality holds (not just $\ge$).  Hence the inequality is
actually an equality in the centred case, and the original $\ge$ becomes $=$.
This is consistent: finite free convolution adds second cumulants.
\end{remark}

%------------------------------------------------------------------
\section*{References}

\begin{itemize}
\item A.~Marcus, D.~Spielman, and N.~Srivastava, ``Finite free convolutions of
  polynomials,'' \emph{Probability Theory and Related Fields}, 182:807--848,
  2022.
\item A.~Marcus, ``Polynomial convolutions and (finite) free probability,''
  \emph{PhD Thesis}, Yale University, 2014.
\item D.~Voiculescu, ``Limit laws for random matrices and free products,''
  \emph{Inventiones Mathematicae}, 104:201--220, 1991.
\end{itemize}

\end{document}
